% ============================================================
% 00_CAPA.TEX
% Sistema Adaptativo de Classificação Multicritério
% de Fluidos de Corte
%
% FINALIDADE
%
% Capa da apresentação.
%
% Este arquivo constitui apenas a BASE.
%
% Informações institucionais, nome da disciplina, nome da
% professora, logos e demais elementos oficiais deverão ser
% preenchidos pelos integrantes quando forem confirmados.
%
% Não inventar informações institucionais.
%
% ============================================================


% ============================================================
% FRAME — CAPA
% ============================================================

\begin{frame}[plain]

    \begin{center}

        % ----------------------------------------------------
        % LOGOS
        % ----------------------------------------------------
        %
        % Quando as logos oficiais estiverem disponíveis,
        % poderão ser inseridas aqui.
        %
        % Exemplo estrutural:
        %
        % \includegraphics[height=1.2cm]{assets/unesp.jpg}
        % \hfill
        % \includegraphics[height=1.2cm]{assets/estat.png}
        %
        % Não ativar enquanto os arquivos reais não existirem.
        %

        \vspace{0.5cm}


        % ----------------------------------------------------
        % INSTITUIÇÃO
        % ----------------------------------------------------

        {\large
        \textbf{INSTITUIÇÃO A CONFIRMAR}
        }

        \vspace{0.15cm}

        {\normalsize
        Curso / Departamento:
        \texttt{TO\_BE\_FILLED}
        }

        \vspace{0.15cm}

        {\normalsize
        Disciplina:
        \texttt{TO\_BE\_FILLED}
        }


        \vfill


        % ----------------------------------------------------
        % TÍTULO
        % ----------------------------------------------------

        {\LARGE\bfseries
        Sistema Adaptativo de Classificação Multicritério
        de Fluidos de Corte
        }

        \vspace{0.5cm}

        {\large
        Auditoria de Dados, CRITIC, TOPSIS, Pareto e
        Análise de Robustez
        }


        \vfill


        % ----------------------------------------------------
        % AUTORES
        % ----------------------------------------------------

        {\large
        \textbf{Integrantes}
        }

        \vspace{0.3cm}

        {\normalsize
        Pacheco
        \qquad
        Benjamin
        \qquad
        Vitor
        \qquad
        Marco
        }

        \vspace{0.6cm}


        % ----------------------------------------------------
        % DOCENTE
        % ----------------------------------------------------

        {\normalsize
        \textbf{Docente:}
        \texttt{TO\_BE\_FILLED}
        }

        \vspace{0.5cm}


        % ----------------------------------------------------
        % DATA / ANO
        % ----------------------------------------------------

        {\normalsize
        2026
        }

        \vspace{0.3cm}

    \end{center}

\end{frame}


% ============================================================
% MATERIAL EXTERNO DA PROFESSORA
% ============================================================
%
% Caso a professora forneça:
%
% - logos;
% - modelo visual;
% - nome oficial da disciplina;
% - nome oficial do trabalho;
% - identificação institucional;
% - orientação específica para a capa;
%
% essas informações deverão substituir os placeholders acima.
%
% Não utilizar logos obtidas aleatoriamente na internet como
% substitutas de material institucional fornecido ou confirmado.
%
% ============================================================


% ============================================================
% LOGOS
% ============================================================
%
% Pasta prevista para os elementos visuais externos:
%
% templates/presentation/latex/assets/
%
% Exemplos futuros:
%
% assets/unesp.jpg
% assets/estat.png
%
% Os nomes definitivos dependerão dos arquivos efetivamente
% recebidos.
%
% ============================================================


% ============================================================
% NOMES DOS INTEGRANTES
% ============================================================
%
% Nesta base utilizamos apenas:
%
% Pacheco
% Benjamin
% Vitor
% Marco
%
% Antes da apresentação final, substituir pelos nomes completos
% conforme a forma oficial de identificação definida pelo grupo
% ou pela professora.
%
% ============================================================


% ============================================================
% TÍTULO
% ============================================================
%
% Título-base do projeto:
%
% Sistema Adaptativo de Classificação Multicritério
% de Fluidos de Corte
%
% Caso a professora determine posteriormente um título oficial,
% utilizar a versão determinada por ela.
%
% ============================================================


% ============================================================
% SUBTÍTULO
% ============================================================
%
% Subtítulo-base:
%
% Auditoria de Dados, CRITIC, TOPSIS, Pareto e
% Análise de Robustez
%
% Ele funciona como descrição metodológica da apresentação e
% poderá ser simplificado pelos integrantes na versão final.
%
% ============================================================


% ============================================================
% ESTADO ATUAL
% ============================================================
%
% COVER_STATUS = BASE_PREPARED
%
% INSTITUTION = TO_BE_FILLED
% COURSE = TO_BE_FILLED
% DEPARTMENT = TO_BE_FILLED
% DISCIPLINE = TO_BE_FILLED
% PROFESSOR = TO_BE_FILLED
% OFFICIAL_LOGOS = TO_BE_FILLED
% FULL_AUTHOR_NAMES = TO_BE_FILLED
%
% ============================================================